\documentclass[12pt]{article}
\usepackage[margin=1in]{geometry}
\usepackage[utf8]{inputenc}
\usepackage[T1]{fontenc}
\usepackage{enumitem}
\usepackage{hyperref}
\usepackage{parskip}
\setlength{\parskip}{6pt plus 2pt minus 1pt}
\usepackage{titlesec}
\usepackage{xcolor}
\usepackage{fancyhdr}
\usepackage{tcolorbox}
\tcbuselibrary{breakable}

% UofSC Color Definitions
\definecolor{garnet}{RGB}{115,0,10}
\definecolor{rose}{RGB}{204,46,64}
\definecolor{atlantic}{RGB}{70,106,159}
\definecolor{congaree}{RGB}{31,65,77}
\definecolor{horseshoe}{RGB}{101,120,11}
\definecolor{neutral90}{RGB}{54,54,54}

% Page setup
\pagestyle{fancy}
\fancyhf{}
\lhead{Daily Gamecock \textit{Opinion}}
\rhead{Writing Nutgrafs}
\cfoot{\thepage}
\renewcommand{\headrulewidth}{0.4pt}

% Section formatting
\titleformat{\section}{\Large\bfseries\color{garnet}}{\thesection}{1em}{}
\titleformat{\subsection}{\large\bfseries\color{neutral90}}{\thesubsection}{1em}{}

% Hyperlink setup
\hypersetup{
    colorlinks=true,
    urlcolor=atlantic,
    linkcolor=atlantic,
    citecolor=atlantic
}

% Define tcolorbox environments
\newtcolorbox{tipbox}[1][]{
    colback=congaree!10,
    colframe=congaree,
    boxrule=1pt,
    sharp corners,
    left=8pt,
    right=8pt,
    top=8pt,
    bottom=8pt,
    title={\textbf{#1}},
    fonttitle=\bfseries,
    coltitle=white,
    colbacktitle=congaree,
    toptitle=4pt,
    bottomtitle=4pt,
    breakable,
    parbox=false
}

\newtcolorbox{promptbox}{
    colback=atlantic!10,
    colframe=atlantic,
    boxrule=1pt,
    sharp corners,
    left=8pt,
    right=8pt,
    top=8pt,
    bottom=8pt,
    breakable,
    parbox=false
}

\newtcolorbox{seasonbox}[1][]{
    colback=horseshoe!10,
    colframe=horseshoe,
    boxrule=1pt,
    sharp corners,
    left=8pt,
    right=8pt,
    top=8pt,
    bottom=8pt,
    title={\textbf{#1}},
    fonttitle=\bfseries,
    coltitle=white,
    colbacktitle=horseshoe,
    toptitle=4pt,
    bottomtitle=4pt,
    breakable,
    parbox=false
}

\newtcolorbox{keybox}{
    colback=garnet!5,
    colframe=garnet,
    boxrule=1pt,
    sharp corners,
    left=8pt,
    right=8pt,
    top=8pt,
    bottom=8pt,
    breakable,
    parbox=false
}

% Custom compact title
\newcommand{\doctitle}{%
    \noindent{\LARGE\bfseries Daily Gamecock \textit{Opinion}}\\[2pt]
    {\LARGE\bfseries Writing Nutgrafs}\\[4pt]
    {\normalsize Opinion Section. Spring 2026}\\[-2pt]
    \noindent\rule{\textwidth}{1pt}
    \vspace{8pt}
}

\begin{document}
\thispagestyle{plain}

\doctitle

\begin{keybox}
A nutgraf is the paragraph that tells your reader what the piece is actually about and why it matters right now. It is the single most important structural element in an opinion piece. This guide will teach you how to write one, how to spot when it's missing, and how to fix it.
\end{keybox}

\section*{What Is a Nutgraf?}

The nutgraf (short for ``nut paragraph'') distills your entire argument into one or two sentences. It answers three questions at once: What is this piece about? Why should the reader care? Why now?

Think of it this way: the nutgraf is the paragraph an editor reads to decide whether your piece is worth running. If your nutgraf is vague, your piece sounds vague. If your nutgraf is sharp, your editor knows exactly what they're getting.

A nutgraf is not the same as a thesis statement from a five-paragraph essay. A thesis states a position. A nutgraf states a position, establishes stakes, and anchors the piece in time. It does more work in fewer words.

\section*{Where Does It Go?}

The nutgraf typically appears in paragraphs 2 through 4 of your piece. The structure looks like this:

\begin{enumerate}[leftmargin=*]
    \item \textbf{Hook / Lede} --- An opening that grabs attention. An anecdote, a striking fact, a scene, a question.
    \item \textbf{Nutgraf} --- The paragraph that tells the reader what this piece is actually arguing and why it matters now.
    \item \textbf{Body} --- Evidence, examples, analysis, counterarguments.
    \item \textbf{Conclusion} --- A closing that ties back to the stakes or offers a call to action.
\end{enumerate}

The reader should never have to wonder ``where is this going?'' past the nutgraf. If someone is still confused about your argument by paragraph 4, the nutgraf is either missing or too weak.

\section*{The Three Questions}

\begin{tipbox}[Every Nutgraf Must Answer These]
\begin{enumerate}[leftmargin=*]
    \item \textbf{What is the argument?} --- State your specific claim. Not ``this is a problem,'' but what you think should happen and who should do it.
    \item \textbf{Why does it matter?} --- What are the stakes? Who is affected? What happens if nothing changes?
    \item \textbf{Why now?} --- What makes this timely? A new policy, a recent event, a deadline, a pattern that just became visible?
\end{enumerate}
If your nutgraf doesn't answer all three, it isn't finished.
\end{tipbox}

\section*{Anatomy of a Strong Nutgraf}

A strong nutgraf has four qualities:

\subsection*{Specificity}
It names the thing. Not ``the university should do better'' but ``the university should fire David Voros.'' Not ``tuition is too high'' but ``USC should freeze tuition increases until state funding returns to 2008 levels.'' The more specific your nutgraf, the harder it is to ignore.

\subsection*{A Clear Stance}
A nutgraf is not a topic sentence. ``College affordability is a growing concern'' is a topic. ``USC's Board of Trustees should cap annual tuition increases at the rate of inflation'' is a stance. Opinion writing requires the latter.

\subsection*{A Sense of Stakes}
The reader needs to know what's at risk. Who gets hurt? What gets worse? What opportunity do we lose? Stakes turn an opinion into an argument worth reading.

\subsection*{A Time Anchor}
Your nutgraf should make clear why you're writing \textit{this week} and not six months ago. A recent vote, a new report, a campus incident, a policy deadline --- something that ties the piece to the present moment.

\section*{Examples from the Archive}

The section's standards have evolved over time. The following examples are drawn from Daily Gamecock opinion pieces published between 2019 and 2021. Each one illustrates a common nutgraf pattern --- some strong, some in need of sharpening. For each piece, we show the original opening, identify the nutgraf (or its absence), and then offer a revised version that meets our current standards.

These are not critiques of the original writers. They are teaching examples. Every writer on staff today will benefit from studying what works and what can be tightened.

%--- Example 1 ---
\subsection*{``Column: Fire David Voros'' --- Audrey Elsberry, 10/20/21}

\begin{promptbox}
\textbf{Original opening:}

``Professor David Voros must be terminated immediately if the university has any ounce of respect for its student body. David Voros is a tenured professor in the university's art department and has been on staff since 2000\ldots\ After three sexual misconduct lawsuits, a petition of more than 3,900 signatures for Voros' termination and the administration's decision to place him on sabbatical rather than formal termination, student safety is being undoubtedly neglected at the hands of the university.''
\end{promptbox}

\textbf{Analysis:} This is one of the stronger examples in the archive. The first sentence is direct and functions as a nutgraf --- it names the person, names the action, and implies the stakes. But the third sentence is actually the strongest nutgraf candidate because it adds the ``why now'': three lawsuits, 3,900 petition signatures, and a sabbatical instead of termination. That accumulation of facts is what makes the argument feel urgent rather than just angry.

\begin{seasonbox}[Revised Nutgraf]
After three sexual misconduct lawsuits, a student petition with more than 3,900 signatures, and an administration that responded by granting a sabbatical rather than pursuing termination, USC must fire art professor David Voros --- or admit that it values tenure over student safety.
\end{seasonbox}

This version compresses the strongest evidence into a single sentence and frames the stakes as a choice the university is making.

%--- Example 2 ---
\subsection*{``Column: College is becoming unaffordable for students'' --- Raymond Escoto, 11/18/21}

\begin{promptbox}
\textbf{Original opening:}

``The cost of college in the U.S. continues to rise to ridiculous heights. This is an unacceptable fact that must change. At the University of South Carolina, the cost of in-state tuition alone (no meals, housing, books or supplies) is \$12,688 for the 2021--2022 school year. Out-of-state tuition is even more inflated at \$33,928. Compare this to USC's tuition in 1977, which, when adjusted for inflation, comes out to \$2,606. There has been a 387\% increase of in-state tuition over the last 44 years.''
\end{promptbox}

\textbf{Analysis:} The first sentence reads like a thesis, but it is too vague --- ``this must change'' doesn't tell us what should change or who should change it. The statistics are strong evidence, but the piece dives into data before stating a specific argument. The reader gets numbers without a frame.

\begin{seasonbox}[Revised Nutgraf]
USC's in-state tuition has risen 387\% since 1977 --- nearly four times the rate of inflation --- and the Board of Trustees has approved an increase every single year. Until the state legislature restores higher education funding to pre-recession levels, the Board should freeze tuition and cut administrative overhead instead of passing costs to students.
\end{seasonbox}

Now the reader knows the argument (freeze tuition, restore state funding), the stakes (students bear the cost), and the time anchor (yearly increases, legislative session).

%--- Example 3 ---
\subsection*{``Column: Recycling can't tackle the plastic pandemic'' --- Christina Ceniccola, 9/30/21}

\begin{promptbox}
\textbf{Original opening:}

``Disposable face masks are creating a serious environmental problem that recycling efforts are not going to solve. In a world that consistently chooses convenience, single-use plastics are the default. These single-use options are increasing the amount of waste we produce worldwide and our environment is feeling the grave effects. The COVID-19 pandemic has introduced a unique facet of plastic pollution that continues to fly under the radar: disposable face masks.''
\end{promptbox}

\textbf{Analysis:} The first sentence is close to a nutgraf, but it's generic --- it could apply to any article about plastic waste. The actual topic (disposable masks and microplastics) doesn't appear until paragraph 4, by which point the reader has waded through broad claims about convenience and waste. The piece needs its sharpest idea pulled up front.

\begin{seasonbox}[Revised Nutgraf]
Billions of disposable face masks entered the waste stream during the COVID-19 pandemic, and unlike most plastics, they cannot be conventionally recycled. As they break down, they release microplastic fibers into waterways and soil --- a pollution crisis that campus recycling programs were never designed to address and that USC has done nothing to mitigate.
\end{seasonbox}

This version names the specific pollutant, explains why recycling fails, and anchors the problem to the pandemic timeline and to USC's inaction.

%--- Example 4 ---
\subsection*{``Column: Students should embrace USC's tradition of protest'' --- Carly McPherson, 11/7/21}

\begin{promptbox}
\textbf{Original opening:}

``In a world consumed with social injustices, humanitarian crises, climate change and global pandemics, we cannot afford the privilege of remaining silent. USC students should embrace and continue the tradition of protesting for their beliefs and follow in the footsteps of former students. Some of us were taught to be timid and diffident, and others to be outspoken and assertive. Ignoring what's comfortable, we should strive to use our voices during these turbulent times.''
\end{promptbox}

\textbf{Analysis:} This is a common pattern --- the nutgraf is too broad and reads like a motivational poster. ``We should use our voices'' doesn't tell the reader what specific argument is coming. There is no ``why now,'' no named event, and no specific call to action. The piece could have been written in any year about any campus.

\begin{seasonbox}[Revised Nutgraf]
Last month, fewer than 200 students showed up to protest the university's handling of the Voros misconduct case --- a turnout that would have been unthinkable during the 1970 anti-war rallies that shut down the Horseshoe. USC has a proud history of student protest, and today's students need to reclaim it before the administration learns it can ignore them entirely.
\end{seasonbox}

Now the reader knows what protest we're talking about, what the historical comparison is, and what the stakes are (administrative indifference).

%--- Example 5 ---
\subsection*{``Editorial: USC administration is unacceptably dishonest'' --- Editorial Board, 11/15/21}

\begin{promptbox}
\textbf{Original opening:}

``The Daily Gamecock has repeatedly witnessed dishonesty and a lack of accountability from university administrators towards students, faculty and staff. While we will continue to report on the administration's actions --- and inactions --- we also must demand that our school's leadership start acting on the promises it makes in the interests of the USC community. We are embarrassed about our own university's lack of accountability and regard for the safety and wellbeing of its students.''
\end{promptbox}

\textbf{Analysis:} The first sentence functions as a nutgraf and it's direct --- the editorial board is calling the administration dishonest. But the word ``repeatedly'' does a lot of work without naming specifics. The piece would be stronger if the nutgraf named the evidence: the FOIA revelations, the building renaming controversy, the gaps in the diversity budget. When you have concrete examples, put them in the nutgraf. That's what makes the accusation credible, not just bold.

\begin{seasonbox}[Revised Nutgraf]
From withholding records requested under FOIA to quietly shelving its own building-renaming recommendations, USC's administration has established a pattern of saying one thing and doing another. The Daily Gamecock's editorial board believes this pattern is not accidental --- it is institutional, and it will continue until students and faculty demand transparency with consequences.
\end{seasonbox}

Naming the specific instances transforms the nutgraf from an emotional claim into an evidentiary one.

%--- Example 6 ---
\subsection*{``Column: USC needs to increase its minimum wage'' --- Varsha Gowda, 10/27/21}

\begin{promptbox}
\textbf{Original opening:}

``The University of South Carolina is underpaying its workers. With the effects of the pandemic in mind, it is crucial that we see an increase to a \$15 minimum wage. Currently, the minimum wage of USC workers is \$7.25 an hour. South Carolina doesn't have a state minimum wage law, meaning the federal standard applies to our state's workers.''
\end{promptbox}

\textbf{Analysis:} This is a decent nutgraf --- it states a clear position (raise wages to \$15) and names the current figure (\$7.25). But ``with the effects of the pandemic in mind'' is doing vague work. What effects? Inflation? Staffing shortages? Workers quitting? A stronger version would name the concrete consequence that makes the raise urgent right now.

\begin{seasonbox}[Revised Nutgraf]
USC pays its lowest-wage workers \$7.25 an hour --- the federal minimum --- while dining halls and maintenance crews remain chronically understaffed because workers can earn more at the Target across the street. The university should raise its minimum wage to \$15 before it loses the workforce that keeps campus running.
\end{seasonbox}

Now the ``why now'' is concrete: the university is losing workers to competitors who pay more. The stakes are operational, not abstract.

\section*{Common Nutgraf Mistakes}

Watch for these patterns in your own drafts:

\begin{itemize}[leftmargin=*]
    \item \textbf{Too vague} --- ``This is a problem that needs to change.'' What problem? What change? Who acts?
    \item \textbf{Too broad} --- ``In today's society\ldots'' This framing adds nothing and delays your argument.
    \item \textbf{Buried} --- The nutgraf appears in paragraph 5 or later. Your reader gave up in paragraph 3.
    \item \textbf{Missing entirely} --- The piece starts arguing without ever framing what the argument is. The reader follows your evidence but never learns where you're headed.
    \item \textbf{Thesis without stakes} --- States an opinion but not why it matters. ``Tuition is too high'' is a thesis. ``Tuition is too high and 40\% of students are taking on debt they won't repay for decades'' has stakes.
    \item \textbf{No time anchor} --- The piece could have been written any year. Nothing ties it to this semester, this news cycle, or this moment on campus.
\end{itemize}

\section*{The Nutgraf Test}

\begin{tipbox}[Self-Check Before You Submit]
Read your nutgraf out loud, in isolation, and ask yourself:

\begin{enumerate}[leftmargin=*]
    \item If an editor read \textbf{only} this paragraph, would they know your \textbf{exact argument}?
    \item Would they understand \textbf{why it matters} --- who is affected and what is at stake?
    \item Would they know \textbf{why you're writing it now} and not six months from now?
\end{enumerate}

If the answer to any of these is no, revise your nutgraf before you revise anything else. Everything in your piece flows from this paragraph.
\end{tipbox}

\section*{Practice Exercise}

For each of the following hypothetical ledes, draft a nutgraf that answers all three questions (argument, stakes, why now). Write 1--2 sentences.

\begin{promptbox}
\textbf{Lede 1:}

``On Monday morning, students arriving at Thomas Cooper Library found every group study room booked solid through finals week --- not by students, but by a campus tutoring company that reserved them in bulk at midnight using an automated system.''

\textit{Your nutgraf:} What should the library do? Who is harmed? Why does this need to be fixed before next finals week?
\end{promptbox}

\begin{promptbox}
\textbf{Lede 2:}

``When the university unveiled its new mental health initiative last spring, administrators promised same-week counseling appointments for any student in distress. Nine months later, the average wait time at the Counseling Center is 17 days.''

\textit{Your nutgraf:} What is the argument --- more funding, more staff, a different model? What happens to students during those 17 days? Why is this urgent right now?
\end{promptbox}

\begin{promptbox}
\textbf{Lede 3:}

``The Student Government Association voted last Thursday to allocate \$40,000 to a speaker series. The vote passed 14--12, with no public comment period and no advance notice to the student body.''

\textit{Your nutgraf:} What's your stance on the spending, the process, or both? What precedent does this set? Why does the timing matter?
\end{promptbox}

\end{document}
